\chapter{Establish the Database Boundary}

The database is our first real service boundary. We will create it before either
Rails application so that later exercises can distinguish database readiness
from application setup problems.

This exercise makes the first persistent change to the jail environment. The
read-only preview comes first, and BSDRunner still requires an exact typed
confirmation before it creates anything.

\section{The proposed service}

We will use the following values throughout the tutorial:

\begin{description}
  \item[starter] PostgreSQL;
  \item[jail name] \texttt{database};
  \item[FreeBSD release] \texttt{15.1-RELEASE};
  \item[private address] \texttt{10.17.89.23/24};
  \item[outbound policy] Internet only;
  \item[starts at boot] disabled; and
  \item[inbound exposure] none.
\end{description}

Disabling autostart is useful while assembling the lab: every lifecycle change
will be deliberate and visible. Internet-only outbound access allows package
installation and updates but does not grant access to private LAN addresses.
PostgreSQL port 5432 is reachable on the managed jail network; it is not
forwarded to the host's LAN interface.

PostgreSQL also needs System V shared memory. BSDRunner gives this jail private
message, semaphore, and shared-memory namespaces by setting
\texttt{sysvmsg}, \texttt{sysvsem}, and \texttt{sysvshm} to FreeBSD's
\texttt{JAIL\_SYS\_NEW} value. Blackship sends these as native jail parameters,
so the TOML manifest encodes that value as the integer~\texttt{1}. Private
namespaces supply the required kernel features without exposing the host's IPC
objects to the jail. The older
\texttt{allow.sysvipc} setting is deliberately not used because it makes the
jail inherit the host namespace.

\section{Exercise 3A: Preview without changing the host}

From the BSDRunner repository, run the same preview used by the graphical
wizard:

\begin{lstlisting}[language=sh]
sh dotfiles/.config/bsdrunner/scripts/\
bsdrunner-jails-backend.sh \
  blackship-preset-preview postgres database \
  15.1-RELEASE 10.17.89.23/24 no internet-only | jq .
\end{lstlisting}

Verify these fields before opening the creation workflow:

\begin{itemize}
  \item \texttt{mutates} is \texttt{false};
  \item \texttt{ready} is \texttt{true};
  \item the selected package is \texttt{postgresql18-server};
  \item the expected private service port is \texttt{5432/tcp};
  \item \texttt{private\_by\_default} is \texttt{true}; and
  \item autostart is \texttt{no}.
\end{itemize}

The listed Blackship operations are the reviewed plan. The preview validates a
temporary Jailfile and manifest and asks Blackship for a dry run. It does not
install PostgreSQL, reserve the address, or create a jail.

\section{Exercise 3B: Create the jail through BSDRunner}

Launch the installed Jails application:

\begin{lstlisting}[language=sh]
sh ~/.config/bsdrunner/scripts/bsdrunner-jails.sh
\end{lstlisting}

Choose \textbf{Create a Service Jail}, then select the \textbf{PostgreSQL}
starter. Walk through the five visible concepts and keep the reviewed values:

\begin{enumerate}
  \item \textbf{Storage}: confirm that the existing Blackship project uses
    \texttt{zroot/bsdrunner/blackship}. Continue without creating another
    project.
  \item \textbf{Userland}: select \texttt{15.1-RELEASE}. Blackship will use an
    isolated jail root for the service. Thick and thin are Bastille creation
    modes, so BSDRunner does not ask for that unrelated choice in a Blackship
    service starter.
  \item \textbf{Network}: keep \texttt{10.17.89.23/24}. The network definition
    already exists; BSDRunner will restore its runtime bridge as part of the
    guarded creation path.
  \item \textbf{Config}: change the generated name to \texttt{database}, keep
    \textbf{Internet only}, and leave \textbf{Start at boot} disabled.
  \item \textbf{Lifecycle}: wait for \textbf{READY TO CREATE}. Compare the
    displayed package, address, policy, and port with Exercise 3A.
\end{enumerate}

Only after the plan matches, type \texttt{database} into the confirmation field
and choose \textbf{Create Jail}. Building the jail may download a FreeBSD
release or packages and can take several minutes.

\section{Capture the one-time credential}

After PostgreSQL starts, BSDRunner creates a database named \texttt{appdb} and a
login role named \texttt{app}. Its generated password is displayed once and is
not persisted in BSDRunner state.

Copy that password into a password manager or another private temporary record.
Do not put it in this repository, the LaTeX source, a screenshot, shell history,
or a commit. If it is lost, the jail detail view offers a guarded
\textbf{Reset database password} action.

The \texttt{appdb}/\texttt{app} pair is a starter smoke-test credential, not the
final application design. A later exercise will create separate databases and
roles for \texttt{store} and \texttt{catalog}, preserving the ownership boundary
defined in Chapter 1.

\section{Checkpoint}

Refresh the Jails view after creation. The new service should be running at
\texttt{10.17.89.23}; its health check may briefly be pending during the
30-second start period before becoming healthy.

The following read-only command provides the same focused runtime view:

\begin{lstlisting}[language=sh]
blackship \
  --config /usr/local/etc/bsdrunner/blackship.toml \
  --format json ps --host | \
jq '.data.jails[] | select(.service_name == "database") | {
  service_name, state, ip, autostart, health
}'
\end{lstlisting}

The checkpoint is complete when:

\begin{itemize}
  \item the service name is \texttt{database};
  \item its state is \texttt{Running};
  \item its address is \texttt{10.17.89.23};
  \item autostart remains \texttt{false};
  \item the PostgreSQL health check becomes healthy; and
  \item BSDRunner shows no inbound port forward for the jail.
\end{itemize}

Stop after this checkpoint. The next exercise will test the database from a
separate jail rather than trusting a successful local service start.
